\vspace{2.5pt plus 3pt minus 0.5pt}\noindent\hspace*{1.2em}\kbvnum{(12)}\hspace{0.55em}ಕರ್ಮಣಿ ಪ್ರಯೋಗದಲ್ಲಿ ಧಾತುವಿನ ಮುಂದೆ ಬರುವ ಪ್ರತ್ಯಯ ಯಾವುದು?\par

\vspace{2.5pt plus 3pt minus 0.5pt}\noindent\hspace*{1.2em}\kbvnum{(13)}\hspace{0.55em}“ಬಾಗಿಲನ್ನು ತೆರೆ,”-ಈ ವಾಕ್ಯವನ್ನು ಕರ್ಮಣಿ ಪ್ರಯೋಗ ದಲ್ಲಿ ಬರೆಯಿರಿ.\par

\vspace{2.5pt plus 3pt minus 0.5pt}\noindent\hspace*{1.2em}\kbvnum{(14)}\hspace{0.55em}“ರಾಮಾಯಣವು ವಾಲ್ಮೀಕಿಯಿಂದ ರಚಿಸಲ್ಪಟ್ಟಿತು.” ಈ ವಾಕ್ಯದ ಕರ್ತರಿ ಪ್ರಯೋಗವನ್ನು ತಿಳಿಸಿ.\par

\vspace{2.5pt plus 3pt minus 0.5pt}\noindent\hspace*{1.2em}\kbvnum{(15)}\hspace{0.55em}“ಬರು” ಧಾತುವಿನ ಭೂತಕಾಲದ ರೂಪಗಳನ್ನು, ಕೆಳಗೆ ಕೊಟ್ಟಿರುವ ವಾಕ್ಯಗಳಲ್ಲಿ ತೆರವಿನ ಎಡೆ ತುಂಬಿರಿ.\par

\vspace{2.5pt plus 3pt minus 0.5pt}\noindent\hspace*{1.2em}\kbvnum{{[}ಅ{]}}\hspace{0.55em}ನೀನೂ ನಾನು ಶಾಲೆಯಿಂದ……\par

\vspace{2.5pt plus 3pt minus 0.5pt}\noindent\hspace*{1.2em}\kbvnum{{[}ಬ{]}}\hspace{0.55em}ನೀನೂ ಅವನೂ ಶಾಲೆಯಿಂದ……\par

\kbvchapter{ಪಾಠ ೧೩}{ಸಮಾಸಗಳು}

\indent “ಅರಸರ ಮನೆ ಸುಂದರವಾಗಿದೆ.”- ಈ ವಾಕ್ಯದಲ್ಲಿ ದಪ್ಪ ಅಕ್ಷರಗಳಲ್ಲಿರುವ ಎರಡು ಪದಗಳಿಗೆ ಬದಲಾಗಿ ಅರಮನೆ ಎಂಬ ಒಂದೇ ಪದವನ್ನು ಬಳಸಬಹುದು. ರೂಢಿಯಲ್ಲಿ ಅರಸರ ಮನೆಯನ್ನು ಅರಮನೆಯೆಂತಲೇ ಕರೆಯುವರು. ಆದುದ ರಿಂದ, “ಅರಸರ ಮನೆ ಸುಂದರವಾಗಿದೆ” ಎಂದು ಹೇಳುವ ಬದಲು “ಅರಮನೆ ಸುಂದರವಾಗಿದೆ” ಎಂದು ಹೇಳಿದರೆ ತಪ್ಪಾಗಲಾರದು. ಅರಸರ, ಮನೆ-ಇವೆರಡು ಪದಗಳನ್ನೂ ಸೇರಿಸಿ ಒಂದೇ ಪದ ಮಾಡಿ ಹೇಳುವುದರಿಂದ ಕಾಲ ಉಳಿತಾಯದ ಜೊತೆಗೆ ಶ್ರಮವೂ ಕಡಿಮೆಯಾಗುವುದು.\par

\indent ಹೀಗೆ, ಎರಡು ಅಥವಾ ಹೆಚ್ಚು ಪದಗಳು ಅರ್ಥ ಬದಲಾವಣೆಯಾಗದ ರೀತಿಯಲ್ಲಿ ಒಂದಾಗಿ ಸೇರಿ ಒಂದು ಪದವಾಗುವುದನ್ನು ಸಮಾಸವೆಂದು ಹೇಳುತ್ತೇವೆ ; ಪದಗಳು ಒಂದಾಗಿ ಸೇರುವಾಗ ಅವುಗಳ ಮಧ್ಯದಲ್ಲಿರುವ ವಿಭಕ್ತಿ ಪ್ರತ್ಯಯ ಲೋಪವಾಗುವುದೆಂಬುದು ಗಮನಿಸಬೇಕಾದ ಅಂಶ.\par

\indent ಉದಾಹರಣೆಗೆ, ತಾವರೆಯ + ಕೊಳ = ತಾವರೆಕೊಳ ಇಲ್ಲಿ “ತಾವರೆಯ” ಎಂಬುದು ಪೂರ್ವಪದವಾಗಿರುವುದು. ಮತ್ತು “ಕೊಳ” ಎಂಬುದು ಉತ್ತರಪದವಾಗಿರುವುದು. ಪೂರ್ವ ಮತ್ತು ಉತ್ತರ\kbvcontinues
